\section{The spectral price of a character table}\label{geo:section}

The lower bound needs a test that sees rank and sparsity at the same
time. To build it, give a matrix two copies of its row index and two
copies of its column index. For a nonzero $N$-by-$N$ matrix $M$ over
the algebraic closure $k$ of the coefficient field, this is the tensor
of the diagonal operator
\[
    C_M(e_x\otimes e_y)=M_{xy}\,e_x\otimes e_y .
\]
Put the input legs at $A,B$ and the output legs at $A',B'$. Bringing
$A$ to $A'$ and $B$ to $B'$ exposes $M$ across the two remaining sites.
Bringing $A$ to $B$ and $A'$ to $B'$ exposes a diagonal matrix with the
entries of $M$ on its diagonal. The two resulting ranks are
$\operatorname{rank}M$ and $|\operatorname{supp}M|$.
Between these placements, the planar price from the companion paper
\cite{oaG} gives a continuous test of the same four-legged tensor.
A character table attains the largest possible price at every rectangle
(Theorem~\ref{geo:saturation}).

\subsection{Prices of placed tensors}

A price gives a lower bound on the cost of supplying a tensor by
wires between sites, allowing linear operations at each site. A
\emph{restriction} is one such operation applied at every site, after
grouping all indices there. On a fixed set of sites, tensor products
group the corresponding factors and direct sums separate the summands
at every site. In Strassen's asymptotic spectrum a \emph{spectral
point} is a restriction-monotone semiring homomorphism from these
tensors to $\mathbb R_{\ge0}$ that preserves zero and one
\cite{StrassenSpectrum}. We use logarithms to base two.

\begin{theorem}[Spectral plane and copying]\label{geo:spectrum}
Over an algebraically closed field $k$ of any characteristic there is an
assignment of a real price $H(T;\mathbf x)$ to each nonzero tensor $T$
with legs at any finite list of positions $\mathbf x$ in the Euclidean
plane. This assignment has the following properties.
\begin{enumerate}
\item A tensor at one site costs zero. For a finite set $S$ of positive
diameter we have $H(T;\mathbf x)=\delta(S)\log\Phi_S(T)$, where $\Phi_S$ is a
spectral point and $\delta(S)$ is half the perimeter of the convex hull
of $S$. For a segment its perimeter is twice its length. The prices
$H$ agree upon insertion of trivial legs.
\item A common label of $q$ values on all sites of a nonempty $S$, represented by
the diagonal tensor $\sum_{a=1}^q\bigotimes_{x\in S}e_a$, costs
$\delta(S)\log q$. In particular a wire of dimension $q$ between $X,Y$
costs $|X-Y|\log q$. Contracting tensor indices placed at a common
site is a local linear map.
\item The price of a fixed tensor is jointly convex and positively
homogeneous in its positions and invariant under simultaneous
translation. At a fixed placement, its sublevel sets, with the zero
tensor adjoined, are Zariski closed in every finite tensor format.
\item Suppose a basis label of width $q$ has the same value on legs at
$X$ and $Y$ throughout the support of $T$. An additional copy of that
label at $Z$ increases the price by at most
\begin{equation}\label{geo:copy}
 \tfrac12\bigl(|X-Z|+|Y-Z|-|X-Y|\bigr)\log q.
\end{equation}
\end{enumerate}
\end{theorem}

For the construction and the copying inequality, see \cite[Theorem~\ext{G}{thm:main}, Section~\ext{G}{sec:spectra} and Propositions~\ext{G}{prop:joint-convexity} and~\ext{G}{lem:submod}]{oaG}. The closedness of spectral sublevels is also proved in \cite{SpectralSublevels}. 
If several pairs of matching labels join
segments through $O$, we may copy all these labels to $O$ at zero
cost. A linear functional there multiplies the coefficients by any
prescribed function of the labels. When the multiplier is nonzero on
the support, the reciprocal gives price equality.

We write $P_M(A,B,A',B')=H(C_M;A,B,A',B')$ for $M\ne0$. To move
a leg, introduce a wire of the same dimension and contract one end
against the leg. The wire price bounds the cost of the move.
For a tensor $T$ put $h_T(\mathbf x)=H(T;\mathbf x)$.
Joint convexity and homogeneity give subadditivity in the positions.
For sums of tensors at a single fixed placement, the appropriate
subadditivity is instead
\begin{equation}\label{geo:sum}
   \Phi_S\Bigl(\sum_j T_j\Bigr)\le\sum_j\Phi_S(T_j).
\end{equation}
Indeed an ordinary sum is a restriction of the direct sum.
Multiplicativity of the spectral point makes the price of a grouped product the sum of the prices of its factors.

\subsection{The rectangle test}

Put the row label $x$ on the left side of a rectangle of height $L$ and
width $W$ and the column label $y$ on the right. The top corners
are $A=(-W/2,L/2)$ and $B=(W/2,L/2)$, with $A',B'$ directly below.
Write $P_M(L,W)$ for the price of this placement, $m=\log N$ and
$D=\sqrt{L^2+W^2}$ for its diagonal length.

Between two sites, row and column operations turn a tensor of flattening
rank $r$ into a wire of dimension $r$. Collapsing a side of the
rectangle therefore gives the rank and support readings
\begin{equation}\label{geo:endpoints}
   P_M(0,W)=W\log\operatorname{rank}M,\qquad
   P_M(L,0)=L\log|\operatorname{supp}M|.
\end{equation}
At the full rectangle the two different labels may both be copied
to its center $O$, as in Figure~\ref{geo:rectangle-figure}.

\begin{figure}[!htbp]
\centering
\begin{tikzpicture}[font=\small,line cap=round,line join=round]
 \begin{scope}[xshift=-4.2cm]
  \node at (0,1.35) {$L=0$};
  \draw[thick,ink] (-.94,0)--(.94,0);
  \fill[sea] (-.94,0) circle (2pt);
  \fill[rust] (.94,0) circle (2pt);
  \node[above] at (-.94,0) {$x$};
  \node[above] at (.94,0) {$y$};
  \node at (0,-.75) {$M$};
 \end{scope}
 \begin{scope}
  \node at (0,1.35) {$L,W>0$};
  \fill[sea!8] (-1.35,.72)--(-1.35,-.72)--(0,0)--cycle;
  \fill[rust!9] (1.35,.72)--(1.35,-.72)--(0,0)--cycle;
  \draw[ink!50] (-1.35,.72)--(1.35,.72) (-1.35,-.72)--(1.35,-.72);
  \draw[thick,sea] (-1.35,.72)--(-1.35,-.72)--(0,0)--cycle;
  \draw[thick,rust] (1.35,.72)--(1.35,-.72)--(0,0)--cycle;
  \node[above left,inner sep=2pt] at (-1.35,.72) {$A$};
  \node[below left,inner sep=2pt] at (-1.35,-.72) {$A'$};
  \node[above right,inner sep=2pt] at (1.35,.72) {$B$};
  \node[below right,inner sep=2pt] at (1.35,-.72) {$B'$};
  \node at (-.93,0) {$x$};
  \node at (.93,0) {$y$};
  \fill[ink] (0,0) circle (1.8pt);
  \node[below,inner sep=2.5pt] at (0,0) {$O$};
 \end{scope}
 \begin{scope}[xshift=4.2cm]
  \node at (0,1.35) {$W=0$};
  \draw[thick,ink] (0,-.72)--(0,.72);
  \fill[ink] (0,.72) circle (2pt) (0,-.72) circle (2pt);
  \node[right] at (0,.72) {$(x,y)$};
  \node[right] at (0,-.72) {$(x,y)$};
 \end{scope}
\end{tikzpicture}
\caption{At height zero the two-site matrix is $M$ in its row and
column labels. At width zero it is diagonal in the pair $(x,y)$.
In the middle, a common row label and a common column label supply the
two shaded triangles. Each triangle has semiperimeter $(L+D)/2$, and
a functional at $O$ multiplies their coefficient by $M_{xy}$.}
\label{geo:rectangle-figure}
\end{figure}

Sharing $x$ among $A,A',O$ and $y$ among $B,B',O$ costs the sum of
the two triangle prices. The local functional
$e_x\otimes e_y\mapsto M_{xy}$ at $O$ supplies the entry for free.
Consequently
\begin{equation}\label{geo:upper}
          P_M(L,W)\le (L+D)m
\end{equation}
even when some of the sites coincide.

\begin{theorem}[Price on a rectangle]\label{geo:saturation}
Let $N\ge1$ and $L,W\ge0$. Every nonzero $N$-by-$N$ matrix over $k$
has rectangle price at most $(L+\sqrt{L^2+W^2})\log N$. For the full
character table $F_G$ of any finite abelian group of order $N$
invertible in $k$,
\[
       P_{F_G}(L,W)=(L+\sqrt{L^2+W^2})\log N .
\]
More generally the supremum of the rectangle prices in any family
satisfying Proposition~\ref{geo:first-phase} equals this upper bound.
\end{theorem}

A permutation matrix, by comparison, uses one common label of size $N$
on the four corners. After relabeling the legs its tensor is precisely
that diagonal, of price $(L+W)m$. At unit width the excess for the
character table is
\[
 P_{F_G}(L,1)-(L+1)m
       = \frac{L^2}{1+\sqrt{1+L^2}}\,m .
\]
This is a constant-order gap when $L$ is of order $1/\sqrt m$.
Section~\ref{bch:proof:section} uses it to distinguish the Fourier
input and output lists from sparse lists. The proof of
Theorem~\ref{geo:saturation} recovers the price at all aspect ratios
from an algebraic operation implemented by two character tables.

\subsection{Why two character tables force the full price}

A \emph{weighted swap} sends $e_x\otimes e_y$ to
$w_{xy}e_y\otimes e_x$, with all $w_{xy}\ne0$. Place its four legs
at $(p,r,-r,-p)$ for any unit vectors $p,r$. An ordinary swap is
the product of two wires on diameters, one carrying $x$ from $p$ to
$-p$ and the other carrying $y$ from $r$ to $-r$. We can copy both
labels to their intersection at zero extra cost and multiply there by
$w_{xy}$ or by $w_{xy}^{-1}$. Hence every weighted swap has the same
price $4m$ at this placement.

Call $C_M$ a \emph{phase} with matrix $M$. A two-phase construction
applies $C_M$ and then $C_D$ to the same two registers, allowing
separate maps $U\otimes V$ before, between and after them.

\begin{lemma}[A swap from two character tables]\label{geo:swap}
Let $G$ be a finite abelian group of order $N$ invertible in $k$,
with all $N$ characters defined over $k$. A two-phase construction
with matrix $F_G$ at both phases realizes a weighted swap.
\end{lemma}

\begin{proof}
Choose a symmetric nondegenerate bicharacter on $G$, taking the
standard pairing on each cyclic factor, to identify the row and
column indices of $F_G$. Permutations of either index act separately
on its two copies. In a character basis every translation
$S_x e_y=e_{y+x}$ is diagonal, with eigenvalues given by the
characters at $x$ up to inversion. Thus $C_{F_G}$, conjugated on the
target register and with the character labels ordered appropriately,
implements the controlled addition $(x,y)\mapsto(x,y+x)$. By symmetry
of the bicharacter the second phase can use the other register as its
control. Negating that control in the phase gives the subtraction
$(u,v)\mapsto(u-v,v)$. After a negation of the first output we have
\[
       (x,y)\ \longmapsto\ (x,y+x)\
         \longmapsto\ (-y,x+y)\ \longmapsto\ (y,x+y).
\]
The last pair differs from a swap by the translation $S_y$ on its
second register. In operator form the composite is $UP$, where $P$
swaps the registers and $U(e_y\otimes z)=e_y\otimes S_yz$. Choose a
character basis matrix $Q$ that diagonalizes all translations
simultaneously. Moving an input basis change through $P$ gives
\[
 (I\otimes Q^{-1})UP(Q\otimes I)
       =\bigl[(I\otimes Q^{-1})U(I\otimes Q)\bigr]P .
\]
The bracketed operator is diagonal and invertible. The two extra maps
on the left act separately on the registers, as required.
\end{proof}

On the diameters a swap uses the entire price $4m$ available for
transporting the four legs from the origin. This equality constrains
how much of the price the first phase must already carry. The next
statement also applies to two-phase constructions with a weighted swap
in their border.

\begin{proposition}[Price forced by a swap]\label{geo:first-phase}
Let $\mathcal A$ be a family of nonzero $N$-by-$N$ matrices over $k$,
with $N\ge2$ and $m=\log N$. Suppose a fixed weighted swap belongs
to the Zariski closure of the two-phase constructions whose first
matrix lies in $\mathcal A$ and whose second matrix is arbitrary.
At any positions $A,B,A',B'$ in the plane,
\begin{equation}\label{geo:four-legs}
 \sup_{M\in\mathcal A}P_M(A,B,A',B')
 \ \ge\ \frac m2
  \bigl(|A-A'|+|B-B'|+|A-B'|+|B-A'|\bigr).
\end{equation}
\end{proposition}

\begin{proof}
A tuple of vectors $g_i$ is a \emph{force} on a nonzero tensor $T$ if it
gives a linear lower bound on the price at every placement, namely
$\sum_i g_i\cdot x_i\le h_T(\mathbf x)$. Since $h_T$ is sublinear
and continuous, it is the support function of its convex set of
forces. Moving leg $i$ bounds the price change by its displacement
times the logarithm of its dimension. Thus $|g_i|$ is at most this
logarithm and the force set is compact. Translation invariance gives
$\sum_i g_i=0$.

Two facts about forces describe the constraints at intermediate sites.
First, a force on the output of a contraction lifts to forces on
its constituents, cancelling in pairs along the contracted wires.
Indeed, give the two ends of each wire the same arbitrary position.
The sum of the constituent prices bounds the output price and is
the support function of the image of the product of their force sets
under addition at the shared sites. The output force, extended by
zero at each internal wire, belongs to this compact convex image by
separation.

Second, if a three-legged tensor has tensor rank at most $q$, the
successive partial sums of any of its forces lie in a disk of diameter
$\log q$. Such a tensor is a restriction of a common $q$-valued
label, so its price is bounded by half the triangle perimeter times
$\log q$. For any cyclic ordering of its three positions $x_i$, this
upper bound is the support function of the triples
$(u_i-u_{i-1})_{i\in\mathbb Z/3}$ with $|u_i|\le\tfrac12\log q$.
Indeed, maximize each $u_i\cdot(x_i-x_{i+1})$ separately. Every
force triple has this form by separation, and its partial sums are
translates of the points $u_i$. When two legs of a tensor are grouped
at one site, the sum of their forces gives a force on the grouped leg
by the defining inequality.

Fix unit vectors $p,r$ and place the external legs at $(p,r,-r,-p)$.
There are two-phase constructions with prices arbitrarily close to $4m$
at this placement. Otherwise a closed price sublevel below $4m$
would contain the weighted swap. Supplying all four legs from the
origin gives the matching upper bound $4m$ on every construction.
Take a sequence whose prices tend to $4m$ and divide all forces by
$m$. Pick a maximizing boundary force for each construction. Its
four components are in the unit disk and their inner products with
$p,r,-r,-p$ sum to a quantity tending to four. The components
therefore tend to $p,r,-r,-p$, respectively.

Lift these forces through the contraction. The separate maps on
the registers may be absorbed on the legs of the phases. By
restriction monotonicity, a force on a phase with these maps applied
is also a force on the original phase. Along a subsequence the two
bounded internal forces have limits. With the legs in input-output
order, the limiting tuples on the first and second phases are
\[
       (p,r,-a,-b),\qquad (a,b,-r,-p),\qquad a+b=p+r .
\]
The last identity is the zero-sum property on the first phase.
To locate $a,b$, group the two column legs of the first phase.
The remaining three-legged tensor has rank at most $N$. Its term
for row $x$ has the single vector $\sum_y M_{xy}e_y^*\otimes e_y$
at the grouped site and one $x$-vector at each other site.
Use the cyclic order $A,A',(B,B')$ in the disk bound. Passing to the
limit puts $0,p,p-a$ in a disk of diameter one. Since $|p|=1$ its
center must be $p/2$, so $a$ is within $1/2$ of $p/2$.

In the second phase group the two row legs and use the order
$B',B,(A,A')$. Now $0,-p,b-p$ lie in a disk with center $-p/2$
and radius $1/2$, placing $b$ in the same disk as $a$. Their midpoint
$(p+r)/2$ is on its boundary because $|r|=1$. Strict convexity forces
$a=b=(p+r)/2$ (Figure~\ref{geo:force-figure}).

\begin{figure}[!htbp]
\centering
\begin{tikzpicture}[font=\small,line cap=round,line join=round]
 \begin{scope}[xshift=-2.2cm]
  \node[draw=sea,fill=sea!6,minimum width=2.3cm,minimum height=5mm] (M) at (0,.8) {$C_M$};
  \node[draw=rust,fill=rust!6,minimum width=2.3cm,minimum height=5mm] (D) at (0,-.8) {$C_D$};
  \draw[ink] (-.62,1.33)--(-.62,1.05) (.62,1.33)--(.62,1.05);
  \node[above] at (-.62,1.29) {$p$};
  \node[above] at (.62,1.29) {$r$};
  \draw[thick,ink] (-.62,.55)--(-.62,-.55) (.62,.55)--(.62,-.55);
  \node[left] at (-.62,.28) {$-a$};
  \node[right] at (.62,.28) {$-b$};
  \node[left] at (-.62,-.28) {$a$};
  \node[right] at (.62,-.28) {$b$};
  \draw[ink] (-.62,-1.05)--(-.62,-1.33) (.62,-1.05)--(.62,-1.33);
  \node[below] at (-.62,-1.29) {$-r$};
  \node[below] at (.62,-1.29) {$-p$};
 \end{scope}
 \begin{scope}[xshift=1.0cm]
  \fill[sea!7] (1,0) circle (1cm);
  \draw[thick,sea] (1,0) circle (1cm);
  \draw[densely dashed,ink!65] (0,0)--(2,0);
  \draw[-{Stealth[length=1.5mm]},thick,ink] (1,0)--(1.5,.866);
  \node[left] at (1.25,.50) {$r/2$};
  \fill[ink] (0,0) circle (1.6pt) (1,0) circle (1.6pt) (2,0) circle (1.6pt);
  \node[left] at (0,0) {$0$};
  \node[below] at (1,0) {$p/2$};
  \node[right] at (2,0) {$p$};
  \draw[rust] (.95,1.184)--(2.05,.548);
  \fill[rust] (1.5,.866) circle (2pt);
  \node[above right,inner sep=2.5pt] at (1.48,.92) {$\displaystyle\frac{a+b}{2}$};
 \end{scope}
\end{tikzpicture}
\caption{A boundary force on the two-phase contraction lifts with opposite
signs on each internal wire (left). The three-leg rank bounds put both
internal vectors $a,b$ in the disk with diameter $[0,p]$ (right).
Their midpoint is $p/2+r/2$ on the supporting tangent, so both vectors
coincide with it.}
\label{geo:force-figure}
\end{figure}

It remains to turn this constraint on the intermediate legs into
distances at an arbitrary placement. Let $M_j\in\mathcal A$ be the
first matrices on the chosen subsequence. Move the input copies of $x,y$
out to $Rp,Rr$ while leaving the output copies at $A',B'$. For each
fixed $R>0$ the limiting first-phase forces give
\[
 \liminf_{j\to\infty}\frac{P_{M_j}(Rp,Rr,A',B')}{m}
       \ \ge\ 2R-\tfrac12(p+r)\cdot(A'+B').
\]
We can implement the move by copying each matched pair to the new
input site and then erasing its old input copy with an all-ones
covector. By \eqref{geo:copy}, the increase in price divided by $m$
is at most
\[
 2R-\tfrac12\bigl(p\cdot(A+A')+r\cdot(B+B')+
                         |A-A'|+|B-B'|\bigr)+o(1),
\]
using $|Rp-X|=R-p\cdot X+o(1)$ as $R\to\infty$. Subtracting the
cost of the move and letting $R\to\infty$ gives
\[
 \sup_{M\in\mathcal A}\frac{2P_M(A,B,A',B')}{m}
 \ \ge\ |A-A'|+|B-B'|+p\cdot(A-B')+r\cdot(B-A').
\]
The first two lengths come from the two pairs of matching labels in
the copying inequality. Maximizing over the independent choices of
$p,r$ gives the two crossed lengths and proves \eqref{geo:four-legs}.
The first matrices used to bound the supremum may depend on $p,r$.
\end{proof}

\begin{proof}[Proof of Theorem~\ref{geo:saturation}]
The case $N=1$ has price zero at every placement. For $N\ge2$ the
upper bound is \eqref{geo:upper}. In Proposition~\ref{geo:first-phase}
the two same-label distances on a rectangle are $L$ and the two
crossed distances are $D$. Its lower bound is therefore $m(L+D)$.
Lemma~\ref{geo:swap} applies with the singleton $\mathcal A=\{F_G\}$.
\end{proof}

\begin{remark}[Convex quadrilaterals]\label{geo:quadrilateral}
The same price computation determines $P_{F_G}$ on four sites in
convex position. If $d_1,d_2$ are the lengths of their hull diagonals,
then
\[
 P_{F_G}(A,B,A',B')
   =\tfrac m2\bigl(|A-A'|+|B-B'|+d_1+d_2\bigr).
\]
The intersection $O$ of the hull diagonals gives an upper bound by
supplying the common labels on $A,A',O$ and $B,B',O$. The four
distances from the vertices to $O$ sum to $d_1+d_2$.
If both diagonals join unlike labels, exchanging the two copies of
one label if necessary places their lengths in the last two terms
of \eqref{geo:four-legs}. Exchanging these copies preserves the tensor
under the coordinate identification of the basis with its dual.
If the diagonals join like labels, their two wires give the tensor
up to nonzero character weights, and both wires pass through $O$.
The copying theorem then preserves their full price under reweighting.
\end{remark}

